\documentclass[12pt]{article}

\usepackage{sbc-template}
\usepackage[utf8]{inputenc}
\usepackage[T1]{fontenc}
\usepackage[brazil]{babel}
\usepackage{graphicx,url}
\usepackage{amsmath,amssymb}
\usepackage{booktabs}
\usepackage{tabularx}
\usepackage{array}
\usepackage{xcolor}
\usepackage{tikz}
\usetikzlibrary{positioning,arrows.meta,fit}

\sloppy

% Marcador para conteúdo que depende da implementação/experimentos
\newcommand{\pend}{\textcolor{red}{\textit{[a preencher]}}}
\newcolumntype{Y}{>{\raggedright\arraybackslash}X}

\title{Mitigação Adaptativa de Ataques DDoS em Redes Definidas\\
  por Software com Aprendizado por Reforço}

\author{Bernardo Euler Barbosa Santos (567068)\inst{1}\\ Renato Rogério Barros Martins (561497)\inst{1}}

\address{Universidade Federal do Ceará (UFC) -- Campus de Russas -- Russas, CE -- Brasil}

\begin{document}

\maketitle

\begin{abstract}
  In Software-Defined Networks (SDN), static DDoS defenses either block
  legitimate users or react too slowly. This paper proposes a deep
  reinforcement learning agent that, for each traffic source, chooses among
  four graduated actions: allow, rate-limit (OpenFlow meters), redirect to an
  analysis host, or temporarily block. The reward penalizes collateral damage
  to legitimate users and the cost of each action. The agent was implemented
  on Mininet, Open vSwitch and os-ken (a maintained Ryu fork), trained in a
  simulator calibrated with emulation traces, and compared with fixed
  thresholds, entropy detection and Random Forest classifiers.
  Preliminary results (one seed) show that the agent reacts fastest (0.3~s) at negligible controller cost and, mostly by redirecting attackers, keeps availability above 99\% under UDP and ICMP floods; under SYN floods, however, it restricts many legitimate sources and trails fixed thresholds and Random Forest.
\end{abstract}

\begin{resumo}
  Em Redes Definidas por Software (SDN), defesas estáticas contra ataques DDoS
  ou bloqueiam usuários legítimos ou reagem tarde demais. Este artigo propõe um
  agente de aprendizado por reforço profundo que, para cada origem de tráfego,
  escolhe entre quatro ações graduadas: permitir, limitar a taxa
  (\textit{meters} OpenFlow), redirecionar para um nó de análise ou bloquear
  temporariamente. A recompensa penaliza o dano colateral a usuários legítimos
  e o custo de cada ação. O agente foi implementado sobre Mininet, Open vSwitch
  e os-ken (\textit{fork} mantido do Ryu), treinado em um simulador calibrado
  com traços da emulação e comparado a limiares fixos, detecção por entropia e
  classificadores Random Forest. Em 112 rodadas no Mininet (duas sementes, cenários S0--S3), o agente reage mais rápido (0,5~s), com custo desprezível no controlador, e, sobretudo redirecionando os atacantes, mantém a disponibilidade acima de 99\% sob inundações UDP e ICMP; sob inundação SYN, porém, restringe muitas origens legítimas e fica atrás do limiar fixo e do Random Forest, com diferença estatisticamente significativa.
\end{resumo}

% =====================================================================
\section{Introdução}\label{sec:intro}

Ataques distribuídos de negação de serviço (\textit{Distributed Denial of
Service} -- DDoS) buscam esgotar recursos de rede ou de processamento de um
alvo por meio de grande volume de tráfego originado de múltiplas fontes,
tornando serviços indisponíveis a usuários legítimos
\cite{mirkovic2004taxonomy}. Apesar de mais de duas décadas de pesquisa, a
mitigação em tempo real permanece difícil porque o tráfego de ataque se mistura
ao tráfego legítimo \cite{liu2018drl}.

As Redes Definidas por Software (\textit{Software-Defined Networking} -- SDN)
separam o plano de controle do plano de dados e oferecem ao controlador uma
visão global da rede \cite{mckeown2008openflow,kreutz2015sdn}. Essa
característica permite coletar estatísticas de fluxos e instalar regras de
descarte, limitação de taxa ou redirecionamento de forma programática. Por outro
lado, a própria SDN introduz novas superfícies de ataque, como a saturação do
controlador por mensagens \textit{packet-in} e o esgotamento das tabelas de
fluxo dos \textit{switches}
\cite{yan2016sdnddos,albuquerque2023slowtcam}.

\textbf{Problema.} Mecanismos de defesa baseados em regras fixas enfrentam um
dilema: um limiar rígido bloqueia usuários legítimos em picos de demanda
(\textit{flash crowds}), enquanto um limiar permissivo ou uma reação lenta
permite a continuidade do ataque. Grande parte da literatura recente concentra-se
na \emph{detecção} supervisionada, frequentemente com classificadores como
Random Forest \cite{bahashwan2023slr,silva2024sbseg,dias2023microsservicos}, e
trata a \emph{resposta} de forma secundária, muitas vezes como uma decisão
binária de bloquear ou não.

\textbf{Quem é afetado.} Usuários legítimos do serviço atacado, que sofrem
indisponibilidade ou bloqueio indevido; operadores de rede, que precisam
calibrar regras manualmente; e o próprio controlador SDN, que pode ser
sobrecarregado durante o ataque.

\textbf{Escopo.} Considera-se uma rede SDN de um único domínio, com um
controlador e topologia em árvore emulada no Mininet, protegendo um servidor
HTTP contra ataques volumétricos (\textit{SYN}, UDP e ICMP \textit{flood}),
inclusive pulsados e multivetor. Ataques à camada de aplicação, ataques com
falsificação de endereços de origem e cenários com múltiplos controladores
ficam fora do escopo.

\textbf{Objetivo geral.} Desenvolver e avaliar um agente de aprendizado por
reforço (\textit{Reinforcement Learning} -- RL) que observe o tráfego em uma
rede SDN e aprenda, para cada origem de tráfego, quando \emph{permitir},
\emph{limitar a taxa}, \emph{redirecionar para análise} ou \emph{bloquear
temporariamente}, reduzindo o ataque sem bloquear usuários legítimos.

\textbf{Objetivos específicos.}
\begin{enumerate}
  \item Construir um ambiente emulado reprodutível com Mininet, Open vSwitch e
        controlador os-ken (\textit{fork} mantido do Ryu), com geração
        automatizada de tráfego legítimo e de ataque.
  \item Formular a mitigação como um Processo de Decisão de Markov (MDP), com
        estados, ações graduadas e recompensa que penaliza dano colateral.
  \item Treinar um agente \textit{Double DQN} e integrá-lo ao controlador.
  \item Comparar o agente com regras de limiar fixo, detecção por entropia e
        detectores supervisionados, sob as mesmas condições, quanto a tráfego de ataque
        bloqueado, falsos positivos, disponibilidade, latência, tempo de
        reação e sobrecarga no controlador.
\end{enumerate}

\textbf{Contribuições.} (i) Uma formulação de mitigação com
\emph{espaço de ações graduado} de quatro níveis, aplicado por origem e com
política compartilhada, que escala com o número de fontes; (ii) uma função de
recompensa que equilibra explicitamente eficácia, dano a usuários legítimos e
custo operacional das ações; (iii) uma avaliação que inclui cenários adversos
não vistos no treinamento -- ataques pulsados, \textit{flash crowds} e enlaces
com atraso e perda --, verificando se a resposta se adapta à intensidade do
ataque e às condições da rede, e que mede a sobrecarga imposta ao controlador;
e (iv) uma implementação aberta e reprodutível -- controlador, simulador
calibrado, agente, \textit{baselines} e \textit{scripts} de experimento --,
disponível no repositório do projeto.

O restante do artigo está organizado da seguinte forma: a
Seção~\ref{sec:fund} apresenta a fundamentação; a Seção~\ref{sec:relacionados}
discute os trabalhos relacionados; a Seção~\ref{sec:metodologia} descreve a
metodologia; a Seção~\ref{sec:arquitetura} detalha a solução proposta; a
Seção~\ref{sec:planejamento} apresenta o planejamento experimental; as
Seções~\ref{sec:resultados} e~\ref{sec:discussao} trazem resultados, discussão e
ameaças à validade; e a Seção~\ref{sec:conclusao} conclui o trabalho.

% =====================================================================
\section{Fundamentação}\label{sec:fund}

\subsection{Redes Definidas por Software e OpenFlow}

Em SDN, a lógica de controle é retirada dos equipamentos de encaminhamento e
centralizada em um controlador, que programa os \textit{switches} por meio de
uma interface padronizada, tipicamente o OpenFlow
\cite{mckeown2008openflow,kreutz2015sdn}. Cada \textit{switch} mantém tabelas
de fluxo cujas entradas possuem campos de casamento (por exemplo, endereços IP,
portas e protocolo), prioridade, contadores e instruções. Na versão 1.3 do
OpenFlow \cite{onf2012openflow13}, destacam-se três recursos usados neste
trabalho:
\begin{itemize}
  \item \textbf{Contadores e estatísticas}: o controlador solicita
        periodicamente estatísticas por fluxo e por porta
        (\textit{FlowStats}/\textit{PortStats}), obtendo pacotes e bytes
        acumulados.
  \item \textbf{Temporizadores}: entradas podem expirar por inatividade
        (\textit{idle\_timeout}) ou após um tempo fixo (\textit{hard\_timeout}),
        o que permite bloqueios temporários que se desfazem automaticamente.
  \item \textbf{Tabela de medidores (\textit{meters})}: um medidor com banda do
        tipo \textit{drop} descarta pacotes que excedem uma taxa configurada,
        implementando limitação de taxa por fluxo ou por origem.
\end{itemize}
O Open vSwitch (OVS) é um \textit{switch} virtual de código aberto com suporte a
OpenFlow \cite{pfaff2015ovs}; o Mininet emula redes inteiras com
\textit{hosts}, \textit{switches} e enlaces em uma única máquina
\cite{lantz2010mininet}; e o Ryu é um arcabouço em Python para construção de
controladores \cite{ryu}.

\subsection{Ataques DDoS}

Ataques DDoS podem ser classificados quanto ao recurso explorado e à dinâmica
da taxa de envio \cite{mirkovic2004taxonomy}. Os ataques de inundação
(\textit{flooding}) volumétricos mais comuns são: \textit{TCP SYN flood}, que
esgota a fila de conexões semiabertas do servidor; \textit{UDP flood} e
\textit{ICMP flood}, que saturam a banda do enlace da vítima. Em ataques de
baixa taxa ou pulsados, o atacante alterna rajadas e silêncio para degradar o
serviço evitando detectores baseados em médias \cite{kuzmanovic2003shrew}. Em
SDN, há ainda ataques ao plano de controle: pacotes com cabeçalhos aleatórios
geram \textit{packet-in} em excesso e sobrecarregam o controlador, ou preenchem
a memória TCAM dos \textit{switches} \cite{yan2016sdnddos,albuquerque2023slowtcam}.

Atributos clássicos para caracterizar tráfego de ataque incluem estatísticas
de fluxo, como pacotes e bytes por fluxo, a proporção de fluxos sem resposta
(assimetria entre ida e volta) e o crescimento do número de portas
\cite{braga2010lightweight}, além da entropia dos endereços de destino, que cai
quando o tráfego converge para uma vítima \cite{mousavi2015entropy}. Também são
comuns a taxa de pacotes por origem, o tamanho médio dos pacotes e a proporção
de segmentos SYN. Bases públicas como a
CIC-DDoS2019 \cite{sharafaldin2019cicddos} e a InSDN \cite{elsayed2020insdn}
são amplamente usadas para avaliar detectores.

\subsection{Aprendizado por Reforço}

No RL, um agente interage com um ambiente modelado como um MDP
$(\mathcal{S},\mathcal{A},P,R,\gamma)$: em cada passo $t$ ele observa o estado
$s_t$, executa a ação $a_t$, recebe a recompensa $r_t$ e o ambiente transita
para $s_{t+1}$ \cite{sutton2018rl}. O objetivo é aprender uma política $\pi$ que
maximize o retorno descontado $G_t=\sum_{k\ge 0}\gamma^k r_{t+k}$. O
\textit{Q-learning} estima a função ação-valor $Q(s,a)$ por diferenças
temporais \cite{watkins1992qlearning}. A \textit{Deep Q-Network} (DQN) aproxima
$Q$ com uma rede neural, utilizando memória de repetição e rede-alvo
\cite{mnih2015dqn}. O \textit{Double DQN} reduz a superestimação dos valores ao
separar a seleção da ação da sua avaliação \cite{vanhasselt2016ddqn}, e a
arquitetura \textit{dueling} decompõe $Q$ em valor do estado e vantagem da ação
\cite{wang2016dueling}. Métodos de gradiente de política, como o PPO
\cite{schulman2017ppo}, são alternativas para espaços de ação maiores.

\subsection{Termos técnicos adotados}

\begin{itemize}
  \item \textbf{Origem}: endereço IP de origem observado na borda da rede; é a
        unidade sobre a qual o agente decide.
  \item \textbf{Intervalo de decisão ($\Delta t$)}: período entre duas coletas
        de estatísticas e, portanto, entre duas decisões do agente.
  \item \textbf{Ação restritiva}: limitar, redirecionar ou bloquear uma origem.
  \item \textbf{Falso positivo}: ação restritiva aplicada a uma origem
        legítima.
  \item \textbf{Dano colateral}: tráfego legítimo perdido, atrasado ou
        rejeitado em consequência das ações de mitigação.
\end{itemize}

% =====================================================================
\section{Trabalhos Relacionados}\label{sec:relacionados}

\textbf{Detecção supervisionada em SDN.} Revisões sistemáticas recentes mostram
que a maior parte dos trabalhos sobre DDoS em SDN trata de detecção com
aprendizado de máquina e aprendizado profundo, com abordagens de classificador
único, híbridas e de \textit{ensemble} \cite{bahashwan2023slr,tebbaa2025slr}. No contexto nacional,
\cite{silva2024sbseg} avaliam cinco algoritmos na base CIC-DDoS2019, com
destaque para o Random Forest, e \cite{dias2023microsservicos} propõem uma
arquitetura de microsserviços para detecção em SDN que permite usar modelos
mais complexos sem prejudicar o desempenho do controlador. Esses trabalhos
alcançam alta acurácia, mas a resposta ao ataque não é o foco, e a decisão
resultante costuma ser binária. Trabalhos anteriores à popularização do
aprendizado profundo já exploravam mapas auto-organizáveis
\cite{braga2010lightweight} e entropia de endereços de destino \cite{mousavi2015entropy}
para detecção leve no controlador.

\textbf{RL para resposta a DDoS.} \cite{malialis2013throttling} propõem o
\textit{Multiagent Router Throttling}, em que roteadores a montante aprendem a
limitar o tráfego destinado à vítima, com extensão posterior voltada à
escalabilidade para um número maior de agentes \cite{malialis2015distributed}.
\cite{liu2018drl} aplicam DRL em SDN para aprender políticas de limitação
contra inundações SYN, UDP e ICMP, superando um método de \textit{router
throttling}. O
\textit{DeepThrottle} usa DRL para ajustar a limitação em SDN, aumentando a
fração de tráfego normal que chega ao servidor em cenários de ataque não vistos
\cite{chen2022deepthrottle}, e \cite{shen2023fast} propõem recompensa com
punição adaptativa e alocação de banda para preservar tráfego normal em
SDN-IoT. \cite{feng2020applayer} tratam ataques na camada de aplicação com uma
recompensa multiobjetivo: o agente mitiga agressivamente quando o servidor está
sobrecarregado, mesmo sacrificando algumas requisições legítimas, e de forma
conservadora nos demais casos, para minimizar o dano colateral -- uma ideia de
resposta proporcional próxima da investigada aqui, porém fora do contexto de
SDN. \cite{yungaicela2022slowrate} combinam um IDS baseado em
aprendizado profundo com um IPS baseado em DRL contra ataques de baixa taxa em
Mininet e ONOS, estendendo a abordagem com NFV e defesa de alvo móvel
\cite{yungaicela2023sdnnfv}. \cite{khozam2024qos} empregam DRL para mitigar
inundações ICMP, SYN e UDP preservando QoS em SDN para 5G, e o
\textit{QoSentry}, dos mesmos autores, utiliza \textit{Double DQN} com
contramedidas adaptáveis, como limitação e redirecionamento de tráfego, e
recompensa orientada à qualidade de serviço percebida por usuários legítimos
(latência e \textit{jitter}) \cite{khozam2025qosentry}. Por fim, o
\textit{MininetGym} oferece um ambiente modular baseado em Mininet para treinar
agentes DQN, PPO e A2C em cenários de segurança \cite{finistrella2025mininetgym}.

A Tabela~\ref{tab:comparativo} sintetiza os trabalhos mais próximos.

\begin{table}[ht]
\centering
\caption{Quadro comparativo dos trabalhos relacionados (n.i.: não informado no
  material consultado; n.a.: não se aplica).}
\label{tab:comparativo}
\scriptsize
\begin{tabularx}{\textwidth}{@{}>{\raggedright\arraybackslash}p{2.4cm}>{\raggedright\arraybackslash}p{1.5cm}YYY>{\raggedright\arraybackslash}p{1.3cm}@{}}
\toprule
\textbf{Trabalho} & \textbf{Técnica} & \textbf{Ações de resposta} &
\textbf{Ambiente / base} & \textbf{Métricas principais} & \textbf{Custo no controlador} \\
\midrule
Malialis e Kudenko (2013; 2015) & RL multiagente & Limitação em roteadores a montante & Simulação (não SDN) & Tráfego legítimo que alcança o servidor & Não \\
Liu et al. (2018) & DRL & Limitação de tráfego & SDN; SYN, UDP e ICMP & Desempenho frente a \textit{router throttling} & n.i. \\
Feng et al. (2020) & RL & Mitigação de requisições (camada 7) & Servidor web (não SDN) & Mensagens maliciosas mitigadas; dano colateral & Não \\
Yungaicela-Naula et al. (2022) & DL + DRL & Mitigação de fluxos de baixa taxa & Mininet + ONOS & Taxa de detecção; sucesso da mitigação & Parcial (amostragem de fluxos) \\
Chen et al. (2022) & DRL & Limitação de tráfego & SDN & Tráfego normal entregue; carga do servidor & n.i. \\
Shen et al. (2023) & DRL & Limitação + alocação de banda & SDN-IoT & Tráfego normal encaminhado; convergência & n.i. \\
Khozam et al. (2024; 2025) & Double DQN & Limitação e redirecionamento & SDN (5G); SYN, UDP e ICMP & QoS de usuários legítimos (latência, \textit{jitter}) & n.i. \\
Silva et al. (2024) & Supervisionado (RF e outros) & Somente detecção & CIC-DDoS2019 & Acurácia, F1 & n.a. \\
Dias et al. (2023) & Supervisionado (RF e outros) & Somente detecção & SDN com microsserviços & F1; desempenho do controlador & Parcial \\
\midrule
\textbf{Este trabalho} & Double DQN & Permitir, limitar (\textit{meter}), redirecionar, bloquear temporariamente & Mininet + OVS + os-ken & Ataque bloqueado, falsos positivos, disponibilidade, latência, reação & Sim \\
\bottomrule
\end{tabularx}
\end{table}

\textbf{Lacuna.} Nos trabalhos analisados, a resposta aprendida é, em geral,
\emph{unidimensional} (um grau de limitação) ou binária (bloquear ou não); a
exceção mais próxima é o \textit{QoSentry} \cite{khozam2025qosentry}, que
combina limitação e redirecionamento. Ainda assim, poucos estudos reúnem, em
uma mesma avaliação, (i) um conjunto de ações de
naturezas diferentes -- limitação, redirecionamento para inspeção e bloqueio
reversível --, (ii) cenários projetados para provocar falsos positivos, como
\textit{flash crowds} legítimos, e (iii) a medição do custo que a própria
política impõe ao controlador. Este trabalho não pretende reproduzir uma
abordagem existente; ele investiga se um espaço de ações graduado, com custos
explícitos por ação, produz respostas proporcionais à intensidade do ataque e
reduz o dano colateral em relação a regras fixas e a detectores
supervisionados.

% =====================================================================
\section{Metodologia}\label{sec:metodologia}

O trabalho segue uma metodologia experimental em etapas, alinhada ao protocolo
proposto para a disciplina:
\begin{enumerate}
  \item \textbf{Delimitação}: definição do problema (mitigação de DDoS
        volumétrico contra um serviço HTTP em uma rede SDN de um único
        domínio), do escopo e das contribuições (Seção~\ref{sec:intro}).
  \item \textbf{Fundamentação e revisão}: estudo de SDN/OpenFlow, ataques DDoS
        e RL, e levantamento de trabalhos por meio de buscadores acadêmicos e
        de metadados (Crossref, Semantic Scholar) cobrindo publicações da IEEE,
        ACM, Elsevier, Springer, AAAI e da SOL/SBC (SBRC, SBSeg, WGRS), com as
        palavras-chave
        \textit{DDoS}, \textit{SDN}, \textit{reinforcement learning},
        \textit{mitigation} e equivalentes em português, priorizando o período
        2018--2026 (Seções~\ref{sec:fund} e~\ref{sec:relacionados}).
  \item \textbf{Planejamento}: formulação de questões de pesquisa, hipóteses,
        \textit{baselines}, métricas e cenários (Seção~\ref{sec:planejamento}).
  \item \textbf{Ambiente e dados}: construção da topologia emulada, coleta
        automatizada de tráfego rotulado, produzido de forma autorizada em
        ambiente isolado, e automação do registro de \textit{logs} e métricas.
        Foram usados apenas dados gerados na própria emulação: bases públicas
        como a CIC-DDoS2019 \cite{sharafaldin2019cicddos} e a InSDN
        \cite{elsayed2020insdn} não foram adotadas porque seus atributos, extraídos
        de capturas completas, não correspondem aos contadores e às amostras
        sFlow disponíveis ao controlador.
  \item \textbf{Implementação-base}: controlador com monitoramento, regras de
        limiar fixo e detector supervisionado; 12 testes unitários (parser
        sFlow, atributos, recompensa, simulador, paridade entre as
        implementações do agente e \textit{baselines}), um teste de fumaça do
        OVS e rodadas curtas de integração.
  \item \textbf{Método proposto}: implementação do ambiente de treinamento, do
        agente e da sua integração com o controlador, com registro de
        hiperparâmetros e auditoria de cada decisão.
  \item \textbf{Experimentos e análise}: execução de todos os métodos nas
        mesmas condições, com repetições, preservação dos dados brutos e testes
        estatísticos.
\end{enumerate}

\textbf{Ambiente computacional.} Os experimentos foram executados em Ubuntu
26.04 sobre WSL2 (núcleo 6.18, Windows 10, 12 núcleos lógicos e 8~GB de RAM),
com Mininet~2.3.0, Open vSwitch~3.7.1 (OpenFlow~1.3, \textit{datapath} do
núcleo), os-ken~4.2.2 -- o \textit{fork} mantido do Ryu \cite{ryu}, de API
compatível, adotado porque o Ryu~4.34 não é mantido desde 2020 e é
incompatível com Python recente --, Python~3.14, PyTorch~2.14, Gymnasium~1.4
\cite{towers2025gymnasium} e scikit-learn~1.9. O núcleo do WSL2 já inclui
como módulos o \texttt{openvswitch} e as disciplinas de fila usadas pelo
Mininet (HTB, \textit{netem}, TBF), dispensando recompilação. Os ataques foram
gerados com \texttt{hping3}, que atinge taxas de milhares de pacotes por segundo
com menor custo de CPU que geradores em Python como o Scapy \cite{scapy}. Todos
os parâmetros e sementes estão em um arquivo de configuração versionado junto
ao código.

% =====================================================================
\section{Arquitetura e Solução Proposta}\label{sec:arquitetura}

\subsection{Visão geral da arquitetura}

A Figura~\ref{fig:arquitetura} apresenta a arquitetura. O plano de dados,
emulado no Mininet, contém clientes legítimos, atacantes, \textit{switches} OVS,
o servidor-vítima (HTTP) e um nó de análise. No plano de controle, a aplicação
os-ken é composta por quatro módulos: (i) \textbf{monitor}, que coleta
\textit{FlowStats}, \textit{PortStats} e amostras sFlow a cada $\Delta t$; (ii)
\textbf{extrator de atributos}, que agrega as estatísticas por origem; (iii)
\textbf{agente}, que escolhe uma ação para cada origem ativa; e (iv)
\textbf{atuador}, que traduz a ação em mensagens OpenFlow. Toda decisão é
registrada com estado observado, valores $Q$ e ação escolhida, permitindo
auditoria \textit{a posteriori}. Para explicar as decisões, calcula-se a
importância de cada atributo por permutação sobre esses registros, medindo
quanto a ação escolhida muda quando o atributo é embaralhado.

\textbf{Coleta de atributos.} Contadores OpenFlow fornecem apenas pacotes e
bytes por regra. Por isso, a tabela~0 de cada \textit{switch} de borda mantém
uma entrada por endereço ativo como origem e outra como destino, que contam o
tráfego e recebem as ações de mitigação, seguidas do encaminhamento na
tabela~1. Atributos que dependem de cabeçalhos (proporção de SYN, portas de
destino distintas e entropia de endereços) são obtidos por amostragem sFlow,
suportada nativamente pelo OVS, com um coletor integrado ao controlador.

\begin{figure}[ht]
\centering
\begin{tikzpicture}[
  font=\footnotesize,
  bloco/.style={draw, rounded corners, align=center, minimum height=0.9cm, text width=2.3cm},
  seta/.style={-{Stealth[length=2mm]}, thick}
]
  % Plano de controle
  \node[bloco] (mon)   at (0,0)   {Monitor de\\estatísticas};
  \node[bloco] (feat)  at (3.2,0) {Extrator de\\atributos};
  \node[bloco] (agent) at (6.4,0) {Agente\\Double DQN};
  \node[bloco] (act)   at (9.6,0) {Atuador de\\políticas};
  \draw[seta] (mon) -- (feat);
  \draw[seta] (feat) -- node[above]{$o_t^i$} (agent);
  \draw[seta] (agent) -- node[above]{$a_t^i$} (act);
  \node[draw, dashed, inner sep=6pt, fit=(mon)(feat)(agent)(act),
        label={[font=\footnotesize\bfseries]above:Controlador os-ken (plano de controle)}] {};

  % Plano de dados
  \node[bloco] (hosts) at (0,-3.4)   {Clientes legítimos\\e atacantes};
  \node[bloco] (ovs)   at (4.8,-3.4) {\textit{Switches} OVS\\(OpenFlow 1.3)};
  \node[bloco] (vit)   at (9.6,-2.8) {Servidor-vítima\\(HTTP)};
  \node[bloco] (anl)   at (9.6,-4.1) {Nó de\\análise};
  \draw[thick] (hosts) -- (ovs);
  \draw[thick] (ovs) -- (vit.west);
  \draw[thick, dashed] (ovs) -- node[below, sloped]{\scriptsize redirecionar} (anl.west);
  \node[draw, dashed, inner sep=6pt, fit=(hosts)(ovs)(vit)(anl),
        label={[font=\footnotesize\bfseries]below:Mininet (plano de dados)}] {};

  % Interação controle-dados
  \draw[seta] (ovs.north west) -- node[left, align=center]{\scriptsize contadores\\\scriptsize OpenFlow e sFlow} (mon.south);
  \draw[seta] (act.south) -- node[right, align=center]{\scriptsize \textit{FlowMod}\\\scriptsize \textit{MeterMod}} (ovs.north east);
\end{tikzpicture}
\caption{Arquitetura da solução proposta.}
\label{fig:arquitetura}
\end{figure}

\subsection{Topologia}

A topologia base é uma árvore de dois níveis: um \textit{switch} de núcleo
conectado ao servidor-vítima, ao nó de análise e a um \textit{host} de
monitoramento, e três \textit{switches} de borda com quatro \textit{hosts}
cada. Entre os 12 \textit{hosts}, a proporção
entre legítimos e atacantes varia por cenário. Os enlaces são configurados com
\texttt{TCLink} (100~Mbit/s, 0{,}5~ms, filas de 100 pacotes), e o enlace da
vítima é o gargalo (10~Mbit/s), para que o ataque produza degradação
mensurável dentro da capacidade de uma única máquina. O servidor HTTP entrega
um arquivo de 20~KB e foi configurado sem \textit{SYN cookies}, com fila de
1.024 conexões semiabertas e duas retransmissões de SYN-ACK, de modo que uma
inundação SYN ocupa a fila por cerca de 7~s. Um dos clientes legítimos, quando
sorteado, faz \textit{upload} contínuo de 2~Mbit/s com \texttt{iperf3},
representando tráfego legítimo de alta taxa. Como os \textit{hosts} são conhecidos, as regras de
encaminhamento e as entradas de contagem por endereço são instaladas de forma
proativa, e uma regra de menor prioridade envia ao controlador apenas o
tráfego de endereços desconhecidos. Assim, o controlador não depende de
\textit{packet-in} a cada novo fluxo.

\subsection{Formulação como MDP}

A decisão é tomada \emph{por origem}, com uma única política compartilhada por
todas as origens ativas. Essa escolha mantém a dimensão do estado fixa,
independentemente do número de \textit{hosts}, e permite reutilizar a
experiência de todas as origens durante o treinamento.

\textbf{Estado.} Para cada origem $i$ no intervalo $t$, a observação é
\begin{equation}
  o_t^i = \left[\, x_t^i \;\|\; g_t \;\|\; \mathrm{onehot}(a_{t-1}^i) \;\|\; \tau_t^i \,\right],
\end{equation}
em que $\|$ denota concatenação; $x_t^i$ reúne atributos locais normalizados (taxa de pacotes e de bytes,
tamanho médio de pacote, proporção de SYN, número estimado de conexões
distintas, número de portas de destino distintas e razão entre tráfego recebido e enviado pela
origem); $g_t$ reúne atributos globais (utilização do enlace da vítima,
latência de sondas enviadas ao servidor por um \textit{host} de monitoramento,
entropia dos IPs de origem calculada pelos contadores de pacotes, taxa de
\textit{packet-in} no controlador e ocupação do nó de análise);
$a_{t-1}^i$ é a ação vigente; e $\tau_t^i$ é o tempo desde a última mudança de
ação. Atributos são calculados sobre janelas deslizantes de alguns intervalos
para capturar ataques pulsados.

\textbf{Ações.} O conjunto $\mathcal{A}=\{a_0,a_1,a_2,a_3\}$ é graduado
(Tabela~\ref{tab:acoes}). O bloqueio usa \textit{hard\_timeout}, de modo que
um erro do agente se desfaz automaticamente; o redirecionamento envia o tráfego
a um nó de análise com capacidade limitada (enlace de 5~Mbit/s), um roteador
Linux que aplica limitação por origem com o módulo \textit{hashlimit} do
\textit{iptables} (100 pacotes/s por origem) e encaminha o tráfego aprovado ao
servidor com latência adicional. Soluções com estado, como o SYNPROXY, foram
descartadas porque exigiriam caminho simétrico, e as respostas do servidor
retornam diretamente aos clientes.

\begin{table}[ht]
\centering
\caption{Ações do agente e sua implementação em OpenFlow 1.3.}
\label{tab:acoes}
\footnotesize
\begin{tabularx}{\textwidth}{@{}p{2.6cm}YY@{}}
\toprule
\textbf{Ação} & \textbf{Implementação} & \textbf{Custo / efeito colateral} \\
\midrule
$a_0$: permitir & Remove regras restritivas da origem; encaminhamento normal & Nenhum, mas o ataque prossegue se a origem for maliciosa \\
$a_1$: limitar a taxa & \textit{FlowMod} com instrução \texttt{meter} e banda \textit{drop} de taxa $r_{\mathrm{lim}}$ & Usuário legítimo limitado sofre perda parcial \\
$a_2$: redirecionar & \textit{FlowMod} que reescreve o MAC de destino e envia o tráfego pela porta do nó de análise & Latência adicional; capacidade de análise limitada \\
$a_3$: bloquear temporariamente & \textit{FlowMod} de descarte com \textit{hard\_timeout} $T_b$ & Usuário legítimo bloqueado fica indisponível por $T_b$ \\
\bottomrule
\end{tabularx}
\end{table}

\textbf{Recompensa.} A recompensa da origem $i$ combina um termo local e um
termo global compartilhado:
\begin{equation}
  r_t^i = \underbrace{w_1\,\mathbf{1}[i \in \mathcal{M}]\,b(a_t^i)
          - w_2\,\mathbf{1}[i \in \mathcal{L}]\,h(a_t^i)}_{\text{local}}
          + \underbrace{w_3\,D_t - w_4\,U^{\mathrm{atq}}_t}_{\text{global}}
          - w_5\,c(a_t^i,a_{t-1}^i),
  \label{eq:recompensa}
\end{equation}
em que $\mathcal{M}$ e $\mathcal{L}$ são os conjuntos de origens maliciosas e
legítimas (conhecidos apenas no ambiente de treinamento); $b$ é o benefício de
aplicar a ação a um atacante, medido pelo efeito real (fração do seu tráfego
efetivamente mitigada no intervalo, de modo que limitar um atacante que já
envia abaixo de $r_{\mathrm{lim}}$ não rende recompensa), e $h$ é o dano a um
usuário legítimo, $h=\max(h_0(a),\,\text{fração do tráfego afetada})$, com piso
$h_0=(0;\,0{,}2;\,0{,}3;\,1)$ para (permitir, limitar, redirecionar, bloquear).
No treinamento, 30\% dos episódios têm atacantes de baixa taxa (3\% a 25\% da
nominal) e 30\% dos clientes legítimos têm taxa até quatro vezes maior, para
criar origens ambíguas.
Separar $b$ e $h$ é necessário: com uma única severidade $\rho(a)$ multiplicando
ambos os termos, o valor esperado local seria linear em $\rho$ e o ótimo seria
sempre uma ação extrema; com $h$ convexo, as ações intermediárias passam a ser
ótimas quando a origem é ambígua (por exemplo, com probabilidade de ataque
0{,}5, limitar ou redirecionar supera permitir e bloquear);
$D_t$ é a disponibilidade do serviço no intervalo (fração de requisições
legítimas atendidas dentro do prazo); $U^{\mathrm{atq}}_t$ é a fração da
capacidade do enlace da vítima ocupada por tráfego de ataque; e $c(\cdot)$
modela o custo operacional de instalar ou trocar regras, desestimulando
oscilações. Com $w_2 > w_1$, um falso positivo custa mais que o benefício de
restringir um atacante, refletindo o objetivo de não bloquear usuários
legítimos. Os rótulos são usados apenas para compor a recompensa no
treinamento; em operação, o agente decide somente a partir de $o_t^i$. Uma
variante sem rótulos, que usa apenas os termos globais observáveis, será
avaliada como ablação.

\subsection{Agente e treinamento}

O agente principal é um \textit{Double DQN} \cite{vanhasselt2016ddqn} com
arquitetura \textit{dueling} \cite{wang2016dueling}: uma rede de duas camadas
ocultas de 128 neurônios (ReLU), otimizador Adam, memória de repetição e
política $\varepsilon$-gulosa, implementado em PyTorch. Como o bloqueio trava
a origem por $T_b$ intervalos, cada transição vai de uma decisão da origem até
a sua próxima decisão, $k$ intervalos depois, com recompensa descontada
acumulada e alvo $r + \gamma^k \max_{a'} Q(s',a')$ (formulação
semi-Markoviana). Para o controlador, os pesos são exportados e a inferência é
feita em \textit{numpy}, sem dependência de bibliotecas de aprendizado
profundo. Os hiperparâmetros utilizados estão na Tabela~\ref{tab:hiper}.

\begin{table}[ht]
\centering
\caption{Hiperparâmetros do agente.}
\label{tab:hiper}
\footnotesize
\begin{tabular}{@{}ll@{}}
\toprule
\textbf{Parâmetro} & \textbf{Valor} \\
\midrule
Intervalo de decisão $\Delta t$ & 1~s \\
Fator de desconto $\gamma$ & 0{,}95 \\
Taxa de aprendizado & $10^{-4}$ \\
Memória de repetição / lote & $10^{5}$ transições / 64 \\
Atualização da rede-alvo & a cada 1.000 passos \\
$\varepsilon$ (exploração) & 1{,}0 $\rightarrow$ 0{,}05 em 100.000 passos \\
Passos de gradiente (simulador) & 300.000 \\
Bloqueio temporário $T_b$ & 10~s \\
Taxa de limitação $r_{\mathrm{lim}}$ & 10\% da capacidade do gargalo (1~Mbit/s) \\
Pesos da recompensa $w_1,\dots,w_5$ & 1; 2; 0{,}5; 0{,}5; 0{,}05 \\
\bottomrule
\end{tabular}
\end{table}

Como cada passo no Mininet ocorre em tempo real ($\Delta t = 1$~s), treinar
exclusivamente na emulação seria lento. Por isso, o treinamento ocorre em duas
fases: (i) um ambiente Gymnasium \cite{towers2025gymnasium} que simula, em
modelo de fluidos, as origens, o enlace gargalo, os medidores e o nó de
análise, com aleatorização de intensidade, número de atacantes e perfis
legítimos; e (ii) ajuste fino com episódios no próprio Mininet, com
exploração $\varepsilon=0{,}1$, reconstruindo as transições a partir dos
registros de decisão do controlador. O simulador é calibrado com execuções sem
defesa (B0) no Mininet: a taxa efetiva do gerador de ataque, os tamanhos de
pacote e as razões de resposta são estimados diretamente dos atributos
observados, e os parâmetros de dinâmica (tamanho e tempo de permanência da fila
de conexões semiabertas e sensibilidade à perda no gargalo) são escolhidos por
busca em grade, minimizando o erro da disponibilidade sem defesa por cenário. Ambas as fases usam apenas os cenários de treinamento
(S0--S3, Tabela~\ref{tab:cenarios}), assim como a escolha dos pesos da
recompensa e dos limiares dos \textit{baselines}. A avaliação final é feita
\emph{somente} no Mininet. O MininetGym \cite{finistrella2025mininetgym} foi
considerado como alternativa de integração, mas não adotado por depender do
controlador OpenDaylight.

\subsection{Baselines}

\begin{itemize}
  \item \textbf{B0 -- Sem defesa}: encaminhamento normal.
  \item \textbf{B1 -- Limiar fixo}: bloqueia temporariamente uma origem cuja
        taxa de pacotes excede um limiar calibrado no cenário normal.
  \item \textbf{B2 -- Entropia}: detecção por queda da entropia, seguida de
        bloqueio temporário das origens cuja participação no tráfego excede
        $1/N$. \cite{mousavi2015entropy} usam a entropia dos endereços de
        destino; como nesta topologia todo o tráfego converge para um único
        servidor, adotou-se a entropia das \emph{origens} do tráfego destinado
        à vítima, com limiar no percentil 1 do cenário normal.
  \item \textbf{B3 -- Supervisionado}: Random Forest \cite{breiman2001rf}
        treinado com os mesmos atributos $x_t^i$ e tráfego rotulado coletado
        nos cenários de treinamento; bloqueia quando a probabilidade de ataque
        excede 0{,}5.
  \item \textbf{B4 -- Supervisionado graduado}: o mesmo Random Forest, com
        faixas de probabilidade mapeadas para as quatro ações. Esse
        \textit{baseline} isola o efeito do RL em relação ao simples uso de
        ações graduadas.
\end{itemize}

% =====================================================================
\section{Planejamento Experimental}\label{sec:planejamento}

\subsection{Questões de pesquisa e hipóteses}

\begin{description}
  \item[QP1 (eficácia).] O agente reduz o tráfego de ataque que alcança a
        vítima com menos falsos positivos que os \textit{baselines} B1--B4?
        \textit{H1: o agente atinge taxa de bloqueio de ataque comparável à de
        B3 com taxa de falsos positivos significativamente menor que B1 e B3.}
  \item[QP2 (adaptação).] A política mantém o desempenho sob padrões de
        tráfego e condições de rede não vistos no treinamento (ataques
        pulsados, \textit{flash crowds}, ataques multivetor e enlaces com
        atraso e perda)? \textit{H2: a degradação de disponibilidade do agente
        nesses cenários é menor que a dos limiares fixos.}
  \item[QP3 (ações graduadas).] Qual é o ganho do espaço de quatro ações em
        relação a um agente com apenas \{permitir, bloquear\}?
        \textit{H3: o espaço graduado reduz o dano colateral sem aumentar o
        tempo de reação.}
  \item[QP4 (custo).] Qual é o tempo de reação e a sobrecarga que o agente
        impõe ao controlador? \textit{H4: o ciclo de decisão (coleta,
        inferência e instalação de regras) permanece abaixo de $\Delta t$, e o
        tempo de reação é de poucos intervalos de decisão, comparável ao de
        B1.}
\end{description}

\subsection{Cenários}

A Tabela~\ref{tab:cenarios} descreve os cenários. O tráfego legítimo é gerado
por clientes HTTP com chegadas de Poisson, tráfego de fundo com
\texttt{iperf3} e sondas ICMP de latência; os ataques são gerados com
\texttt{hping3}. Cada execução dura 180~s: 30~s de aquecimento, 120~s de
ataque e 30~s de recuperação. Cenários marcados como \textit{não vistos} não
são usados no treinamento.

\begin{table}[ht]
\centering
\caption{Cenários experimentais.}
\label{tab:cenarios}
\footnotesize
\begin{tabularx}{\textwidth}{@{}p{0.9cm}p{3.6cm}Yp{1.6cm}@{}}
\toprule
\textbf{Id} & \textbf{Cenário} & \textbf{Descrição} & \textbf{Treino?} \\
\midrule
S0 & Normal & Apenas tráfego legítimo (mede falsos positivos de base) & Sim \\
S1 & \textit{SYN flood} & 3 intensidades (baixa, média, alta), 2 a 4 atacantes & Sim \\
S2 & \textit{UDP flood} & 3 intensidades & Sim \\
S3 & \textit{ICMP flood} & 3 intensidades & Sim \\
S4 & Ataque pulsado & Rajadas \textit{on-off} de UDP/SYN (baixa taxa média) & Não \\
S5 & \textit{Flash crowd} & Aumento abrupto de clientes legítimos, sem ataque & Não \\
S6 & Multivetor & SYN + UDP simultâneos, com mais atacantes que no treino & Não \\
S7 & Rede degradada & \textit{SYN flood} de intensidade média com 20~ms de atraso e 1\% de perda nos enlaces de borda (\texttt{netem}) & Não \\
S8 & Atacantes furtivos & SYN/UDP de baixa taxa (3\% a 25\% da intensidade baixa) e clientes legítimos intensos, gerando origens ambíguas & Não \\
\bottomrule
\end{tabularx}
\end{table}

\subsection{Métricas}

Seja $B^{\mathrm{atq}}_{\mathrm{env}}$ o volume de ataque enviado e
$B^{\mathrm{atq}}_{\mathrm{mit}}$ o volume de ataque descartado pelas regras de
mitigação (medidores, regras de descarte e nó de análise), obtido dos
contadores dessas regras. Descartes causados apenas pelo congestionamento do
gargalo não contam como mitigação. As métricas são:
\begin{itemize}
  \item \textbf{Taxa de ataque bloqueado}:
        $\mathrm{TAB} = B^{\mathrm{atq}}_{\mathrm{mit}} / B^{\mathrm{atq}}_{\mathrm{env}}$.
  \item \textbf{Usuários legítimos bloqueados (falsos positivos)}: fração de
        origens legítimas que receberam ação restritiva em algum intervalo, e
        fração do tempo em que permaneceram restritas, reportando
        separadamente limitação/bloqueio e redirecionamento, já que este último
        apenas acrescenta latência.
  \item \textbf{Disponibilidade}: fração de requisições HTTP legítimas
        concluídas com sucesso em até 2~s.
  \item \textbf{Latência e perda}: tempo de resposta HTTP (mediana e p95),
        RTT das sondas e perda de pacotes legítimos.
  \item \textbf{Tempo de reação}: intervalo entre o início do ataque e a
        primeira ação restritiva efetiva sobre cada atacante; e tempo até a
        disponibilidade voltar a 95\% do valor medido no cenário S0.
  \item \textbf{Sobrecarga no controlador}: uso de CPU e memória do processo
        do controlador, mensagens OpenFlow por segundo (\textit{FlowMod},
        \textit{MeterMod}, estatísticas), tempo de inferência por decisão e
        número de regras instaladas.
\end{itemize}

\subsection{Repetições e controle de variáveis}

Cada combinação de método e cenário foi executada com 3 sementes distintas e
fixadas (o protocolo prevê até 10), reiniciando a rede entre execuções; todos
os métodos usam as mesmas sementes, topologia, perfis de tráfego e intervalo de
coleta. Os resultados são reportados como média e intervalo de confiança de
95\% (\textit{bootstrap}). Para comparar o agente com cada \textit{baseline},
aplica-se o teste de postos sinalizados de Wilcoxon pareado por cenário e
semente, com correção de Holm para comparações múltiplas. Todos os resultados,
inclusive os desfavoráveis, são reportados, e os dados brutos (estatísticas,
registros de decisão e \textit{logs} dos clientes) são preservados. Com seis
métodos, 15 configurações de cenário (S1--S3 em três intensidades cada) e
execuções de 3~min, a campanha principal soma cerca de 14~horas de emulação,
executadas em lote por \textit{scripts}; as ablações foram restritas aos
cenários S1H, S4, S5 e S8.

% =====================================================================
\section{Resultados}\label{sec:resultados}

\textbf{Estado da avaliação.} Os resultados a seguir correspondem a 112
rodadas da campanha principal (180~s por rodada, sementes 1001 e 2001),
cobrindo os cenários S0--S3M para os seis métodos; em S3H, apenas B0--B3
foram executados. Os cenários de adaptação (S4--S8) e as ablações não foram
executados no Mininet. Os intervalos de confiança são calculados por
\textit{bootstrap} sobre as rodadas (cenário $\times$ semente) e, portanto,
refletem sobretudo a variação entre cenários. O agente avaliado passou pelo
ajuste fino no Mininet (12 episódios, $\epsilon=0{,}1$) com a recompensa
baseada em efeito.

\subsection{Linha de base sem defesa e calibração}

Sem defesa (B0, 180~s), a disponibilidade foi de 100\% em S0; 99,9\%, 93,7\% e
7,8\% em S1L/M/H; 99,9\%, 21,1\% e 1,9\% em S2L/M/H; e 99,9\%, 98,8\% e 20,1\%
em S3L/M/H. A calibração com os três tipos de ataque reduziu o erro
quadrático da disponibilidade simulada para 0,0175.

\subsection{Eficácia e dano colateral no Mininet (QP1)}

\begin{table}[ht]
\centering
\caption{Resultados no Mininet, média dos cenários S1L--S3M
  (180~s, duas sementes, 16 rodadas por método; $\pm$ meia-largura do IC
  95\% por \textit{bootstrap} entre rodadas). TAB: tráfego de ataque descartado por limitação ou
  bloqueio (não inclui o redirecionado); Legít.: fração de origens legítimas
  que receberam alguma ação restritiva.}
\label{tab:res-eficacia}
\footnotesize
\begin{tabular}{@{}lccccc@{}}
\toprule
\textbf{Método} & \textbf{TAB (\%)} & \textbf{Legít. (\%)} &
\textbf{Disp. (\%)} & \textbf{Lat. p95 (ms)} & \textbf{Reação (s)} \\
\midrule
B0 -- Sem defesa & 0,0 & 0,0 & 67,3 $\pm$ 18,7 & 885 & -- \\
B1 -- Limiar fixo & 93,9 $\pm$ 0,6 & 0,0 & 91,3 $\pm$ 7,4 & 606 & 1,7 \\
B2 -- Entropia & 93,1 $\pm$ 1,3 & 100,0 & 24,8 $\pm$ 3,2 & 1.091 & 2,7 \\
B3 -- RF & 94,4 $\pm$ 0,6 & 1,5 $\pm$ 1,8 & 91,7 $\pm$ 7,9 & 570 & 1,0 \\
B4 -- RF graduado & 94,5 $\pm$ 0,8 & 0,0 & 91,1 $\pm$ 8,4 & 523 & 1,0 \\
Agente (proposto) & 60,6 $\pm$ 15,6 & 37,3 $\pm$ 16,5 & 71,4 $\pm$ 18,0 & 534 & 0,5 \\
\bottomrule
\end{tabular}
\end{table}

\begin{table}[ht]
\centering
\caption{Agente por cenário no Mininet (média de duas sementes). Tempo restr.: fração do
  tempo de ataque em que as origens legítimas estiveram sob ação restritiva.}
\label{tab:res-agente}
\footnotesize
\begin{tabular}{@{}lccccc@{}}
\toprule
\textbf{Cenário} & \textbf{TAB (\%)} & \textbf{Legít. (\%)} &
\textbf{Tempo restr. (\%)} & \textbf{Disp. agente (\%)} & \textbf{Disp. B0 (\%)} \\
\midrule
S0 (normal)   & --   & 12,5 & 2,8 & 99,1 & 100,0 \\
S1L (SYN)     & 79,8 & 56,2 & 2,6 & 36,3 & 99,9 \\
S1M (SYN)     & 93,1 & 83,8 & 5,6 & 26,0 & 95,6 \\
S1H (SYN)     & 73,0 & 88,8 & 5,4 & 24,7 & 24,4 \\
S2L (UDP)     & 57,5 & 11,2 & 0,8 & 99,6 & 100,0 \\
S2M (UDP)     & 23,1 & 11,2 & 0,8 & 99,6 & 38,5 \\
S2H (UDP)     & 9,7  & 11,2 & 0,8 & 99,6 & 3,7 \\
S3L (ICMP)    & 77,6 & 11,2 & 0,8 & 99,6 & 100,0 \\
S3M (ICMP)    & 81,5 & 12,5 & 0,7 & 99,7 & 99,5 \\
\bottomrule
\end{tabular}
\end{table}

Os \textit{baselines} B1, B3 e B4 descartaram cerca de 93--95\% do tráfego de
ataque quase sem restringir usuários legítimos e elevaram a disponibilidade
média de 67,3\% para cerca de 91\% (Tabela~\ref{tab:res-eficacia}). B2
(entropia) restringiu todas as origens legítimas, efeito do limiar estimado
sobre um tráfego normal de baixa entropia. O agente reagiu mais rápido que
todos (0,5~s), mas, na média, ficou próximo da linha sem defesa. O teste de
Wilcoxon pareado por rodada (15 pares) com correção de Holm indica que a
diferença é significativa: o agente teve disponibilidade menor que B1, B3 e
B4 (diferença média de 19 pontos percentuais, $p_{\text{Holm}} < 0{,}01$),
menor TAB (cerca de 33 pontos) e mais origens legítimas restringidas (cerca
de 36 pontos), e foi superior apenas a B2 (47 pontos de disponibilidade,
$p_{\text{Holm}} < 0{,}001$).

A média esconde dois comportamentos opostos (Tabela~\ref{tab:res-agente}).
Sob UDP e ICMP, o agente manteve 99,6--99,7\% de disponibilidade, inclusive
em S2H, onde sem defesa ela cai para 3,7\%, embora a TAB sob UDP tenha ficado
entre 10\% e 58\%. A diferença se explica pela ação escolhida: em cerca de 95\% dos
intervalos de decisão sobre atacantes, o agente \emph{redirecionou} o tráfego
para o nó de análise, que o retira do caminho da vítima mas não é contado como
descarte pela TAB. Sob SYN, ao contrário, a disponibilidade caiu para
25--36\%, abaixo da linha sem defesa em S1L e S1M. O agente restringiu, em
algum momento, de 56\% a 89\% das origens legítimas, quase sempre com
limitação de taxa e por pouco tempo (3--6\% do período de ataque); essa
restrição curta, sozinha, não parece explicar a queda, cuja causa ainda está
sendo investigada (ver Seção~\ref{sec:discussao}). Em S0, sem ataque, o agente
restringiu 12,5\% das origens legítimas, mantendo 99,1\% de disponibilidade.

\subsection{Adaptação e uso das ações (QP2 e QP3)}

Os cenários de adaptação (S4--S8) ainda não foram executados no Mininet. No
simulador calibrado, o agente obteve 97\% de disponibilidade no ataque
pulsado (S4), contra 93,5\% de B1 e B3, e a variante \{permitir, bloquear\}
chegou a 100\%. A Tabela~\ref{tab:res-acoes} mostra a distribuição das ações
no simulador com origens ambíguas.

\begin{table}[ht]
\centering
\caption{Distribuição das ações durante o ataque no simulador com atacantes
  furtivos (40 episódios).}
\label{tab:res-acoes}
\footnotesize
\begin{tabular}{@{}llcccc@{}}
\toprule
\textbf{Política} & \textbf{Origens} & \textbf{Permitir} & \textbf{Limitar} &
\textbf{Redirecionar} & \textbf{Bloquear} \\
\midrule
Agente (4 ações) & atacantes comuns   & 9,5\% & 2,7\% & 3,6\% & 84,3\% \\
                 & atacantes furtivos & 12,1\% & 0\% & 6,9\% & 81,0\% \\
                 & legítimas          & 100\% & 0\% & 0\% & 0\% \\
Agente \{permitir, bloquear\} & atacantes furtivos & 13,5\% & -- & -- & 86,5\% \\
B4 -- RF graduado & atacantes furtivos & 99,7\% & 0,3\% & 0\% & 0\% \\
\bottomrule
\end{tabular}
\end{table}

No simulador, as ações intermediárias foram pouco usadas (cerca de 6\% das
decisões sobre atacantes). No Mininet, depois do ajuste fino, o quadro se
inverteu: o redirecionamento passou a ser a ação dominante sobre atacantes
(cerca de 95\% das decisões) e a limitação, a única restrição frequente sobre
legítimos. O ajuste fino, portanto, alterou de forma substancial a política,
e as ações graduadas passaram a ser usadas; se esse uso traz vantagem (H3)
depende das ablações, ainda em execução.

\subsection{Custo no controlador (QP4)}

\begin{table}[ht]
\centering
\caption{Custo no controlador no cenário S1H (180~s, média de duas sementes).}
\label{tab:res-custo}
\footnotesize
\begin{tabular}{@{}lcccc@{}}
\toprule
\textbf{Método} & \textbf{CPU (\%)} & \textbf{Memória (MB)} &
\textbf{Msgs OF/s} & \textbf{Decisão (ms)} \\
\midrule
B1 -- Limiar fixo & 1,4 & 190 & 14,9 & 0,17 \\
B3 -- RF          & 4,5 & 192 & 14,9 & 32,11 \\
Agente            & 1,3 & 68  & 15,1 & 0,48 \\
\bottomrule
\end{tabular}
\end{table}

O ciclo de decisão do agente levou cerca de 0,5~ms, três ordens de grandeza
abaixo de $\Delta t=1$~s, com CPU equivalente à do limiar fixo e menos
memória; o Random Forest foi o mais custoso (cerca de 32~ms por ciclo). A
hipótese H4 foi confirmada.

% =====================================================================
\section{Discussão e Ameaças à Validade}\label{sec:discussao}

\textbf{Discussão.} Os resultados dependem do tipo de ataque.
Contra inundações volumétricas (UDP e ICMP), o agente protegeu a vítima tão bem
quanto os \textit{baselines}, reagiu mais rápido e afetou pouco os legítimos,
preferindo redirecionar a bloquear. Contra inundação SYN, porém, foi pior que
os \textit{baselines} e, em dois dos três níveis, pior que não defender (H1 e
H2 não confirmadas; a diferença para B1, B3 e B4 é significativa). Três explicações são plausíveis e serão
verificadas. (i) Redirecionar um SYN \textit{flood} protege menos que
bloqueá-lo: as conexões semiabertas já acumuladas continuam ocupando a fila do
servidor, e a recompensa por efeito trata redirecionamento e bloqueio de modo
parecido. (ii) Lacuna entre simulação e emulação: sob SYN, as conexões
legítimas que não completam o \textit{handshake} alteram as \textit{features}
dos clientes (razão SYN/ACK, pacotes por fluxo) de forma não reproduzida no
simulador, o que leva o agente a limitá-los. (iii) O ajuste fino foi curto
(12 episódios) e com exploração ($\epsilon=0{,}1$), o que pode ter reforçado
escolhas feitas ao acaso. A comparação também expõe uma limitação da TAB, que
subestima a mitigação de quem redireciona; por isso a disponibilidade é a
métrica principal desta análise. Por fim, os \textit{baselines}
supervisionados mostram que as inundações do \texttt{hping3} são facilmente
separáveis, de modo que bloquear rapidamente já é uma resposta muito boa.

\textbf{Validade interna.} O agente foi treinado no simulador e ajustado com
poucos episódios no Mininet. O Mininet compartilha a CPU da máquina hospedeira,
o que pode afetar latência e custo; todos os métodos foram executados nas
mesmas condições. A recompensa usa rótulos disponíveis apenas na emulação.

\textbf{Validade externa.} A topologia é pequena, há um único controlador e o
tráfego é sintético; os resultados podem não se generalizar para redes de
produção, ataques com falsificação de endereços (\textit{IP spoofing}) ou
ataques à camada de aplicação, que estão fora do escopo.

\textbf{Validade de construto.} A disponibilidade é medida por requisições
HTTP com prazo fixo, e o falso positivo é definido no nível de origem: uma
limitação breve conta tanto quanto um bloqueio longo, razão pela qual também se
reporta a fração do tempo sob restrição. A TAB não conta o tráfego
redirecionado.

\textbf{Validade de conclusão.} Há uma única semente e só os cenários S0--S2 e
S3L; nenhuma diferença é estatisticamente significativa após a correção de
Holm. As conclusões são indicativas e dependem da campanha completa, em
execução.

\textbf{Aspectos éticos.} Todo tráfego de ataque é gerado em ambiente
emulado e isolado, sem atingir redes ou serviços reais.

% =====================================================================
\section{Conclusão}\label{sec:conclusao}

Este trabalho propôs e implementou um agente de aprendizado por reforço
profundo (Double Dueling DQN) para mitigação adaptativa de ataques DDoS em SDN,
que escolhe, para cada origem de tráfego, entre permitir, limitar a taxa,
redirecionar para análise ou bloquear temporariamente. A solução completa --
controlador os-ken com contadores por origem e sFlow, simulador calibrado com
traços do Mininet, cinco \textit{baselines}, ajuste fino na emulação e
\textit{scripts} reprodutíveis de experimento e análise -- foi construída sobre
ferramentas abertas. Em 112 rodadas no Mininet, o agente foi o mais rápido
(0,5~s), teve custo desprezível no controlador e, contra inundações UDP e ICMP,
manteve a disponibilidade acima de 99\% usando sobretudo o redirecionamento.
Contra inundação SYN, porém, restringiu muitos usuários legítimos e ficou
atrás do limiar fixo e do Random Forest, que tiveram o melhor desempenho geral.

Como trabalhos futuros, destacam-se: concluir a campanha com três sementes, os
cenários de adaptação (S4--S8) e as ablações; diagnosticar a queda sob SYN,
diferenciando na recompensa o efeito do redirecionamento e do bloqueio sobre
a fila de conexões; modelar no simulador o efeito do SYN sobre as
\textit{features} dos clientes legítimos; fazer um ajuste fino mais longo;
contabilizar o tráfego redirecionado na métrica de mitigação; e estender a
proposta para múltiplos controladores e ataques com falsificação de endereços.

\bibliographystyle{sbc}
\bibliography{referencias}

\end{document}
